\chapter{The incidents list}
\label{ch:incidents}

\shot{05-incidents-list.png}{The incidents list, signed in as an incident
commander. This is the screen people leave open all day.}

\section{Reading the screen}

Seven columns, and each answers a question an on-call engineer actually asks.

\begin{center}
\begin{tabular}{@{}>{\raggedright\arraybackslash}p{4.4cm}p{9.3cm}@{}}
\toprule
\textbf{Column} & \textbf{The question it answers} \\
\midrule
Reference & ``Which incident are we talking about?'' A short stable code you can say out loud on a call. \\
Severity & ``How bad?'' \\
Status & ``Is anyone on it?'' \\
Title & ``What is it?'' \\
Service & ``What is affected?'' \\
Opened & ``When did this start?'' \\
Running / resolved in & ``How long has this been going on --- or how long did it take?'' \\
\bottomrule
\end{tabular}
\end{center}

\note{That last column changes meaning depending on the row. For a live incident
it counts up; for a finished one it is a fixed duration. One column, two
meanings, and it works because you always want the same thing: elapsed time that
matters right now.}

\section{The button that is not always there}

\where{web/src/pages/IncidentsPage.tsx}
\begin{lstlisting}[language=tsx]
{can('incident.create') && (
  <button
    type="button"
    onClick={() => setCreating(true)}
    className="btn btn-primary"
  >
    <Plus className="h-4 w-4" aria-hidden="true" />
    Report incident
  </button>
)}
\end{lstlisting}

That \fpath{\&\&} is JavaScript's short-circuit: if \fpath{can('incident.create')}
is false, the expression evaluates to false and React renders nothing. The
button does not appear disabled --- it does not exist.

Compare the same screen seen by a viewer:

\shot{11-incidents-as-viewer.png}{The identical page, signed in as Sofia Rossi,
a viewer. The ``Report incident'' button is gone. So is ``Export CSV''.}

\warnbox{This is presentation only. A viewer who opens the browser console and
sends \fpath{POST /api/v1/incidents} directly is not stopped by a missing
button, because the button was never the control. The API refuses the request
using a policy on the server. Hiding it is politeness --- it keeps people from
attempting things they cannot do. Chapter~\ref{ch:roles} shows both halves side
by side.}

\section{Where \fpath{can()} comes from}

The frontend is told, at sign-in, exactly what this user may do. It does not
guess from the role name --- it receives a list of permission strings and checks
membership. That matters: the mapping from role to permissions lives in the API,
in one place, and the frontend never has a second copy of it that could drift.

\section{Every badge is an enum}

\where{web/src/pages/IncidentsPage.tsx}
\begin{lstlisting}[language=tsx]
<td className="px-3 py-2">
  <SeverityBadge severity={incident.severity.value} label={incident.severity.label} />
</td>
<td className="px-3 py-2">
  <StatusBadge status={incident.status.value} label={incident.status.label} />
</td>
\end{lstlisting}

Notice the API sends \emph{both} \fpath{value} and \fpath{label}: the machine
form (\fpath{sev1}) and the human form (\fpath{Critical}). The badge colours
itself from \fpath{value} and prints \fpath{label}.

\note{The alternative --- sending only \fpath{sev1} and having the frontend hold
a lookup table of display names --- means two copies of the same knowledge. Add
a severity level and you must remember to update the other side. Sending both
means the API stays the single source of truth for what its own values mean.}

\section{The colour rule this whole design follows}

\where{web/src/index.css}
\begin{lstlisting}
/* red means severity ... So the brand itself is deliberately not red. */
\end{lstlisting}

In an incident tool, red already has a job: it means \emph{this is bad}. If the
brand colour were also red, every page would be permanently alarming and a real
\fpath{sev1} would not stand out. So the brand is indigo, and red is reserved
for severity alone.

That single rule explains the entire palette --- and it is the sort of decision
that is obvious in hindsight and invisible unless someone writes it down.

\section{Filtering, and why the URL changes}

The filter bar above the table --- search, severity, status, service --- writes
its state into the page's URL. That is why you can copy the address of
``all open sev1 incidents on the payments service'' and send it to a colleague,
and why the browser Back button steps through your filters.

\note{Storing view state in the URL rather than in a variable is a small amount
of extra work that turns a screen into something shareable. On an incident call,
pasting a link is faster than describing which dropdowns to set.}

\section{Check your understanding}

\begin{itemize}[leftmargin=1.4cm]
\cbox Why does the API send both \fpath{value} and \fpath{label} for a severity?
\cbox A viewer cannot see ``Report incident''. What actually stops them creating
  an incident?
\cbox Why is the brand colour not red?
\cbox What is gained by putting filter state in the URL?
\end{itemize}
